\section{Related Work}
\label{sec:related}

\paragraph{Provenance Graph Intrusion Detection.}
Provenance-based intrusion detection has evolved from feature-engineering approaches to graph neural network (GNN) models.
Unicorn~\cite{unicorn} extracts histogram-based graph features from provenance graphs, enabling streaming detection without labeled attack data.
ProvDetector~\cite{provdetector} identifies anomalous execution paths by embedding causal paths and flagging outliers.
ThreaTrace~\cite{threatrace} formulates detection as a graph-level classification task, assigning threat labels to entire subgraphs.
ShadeWatcher~\cite{shadewatcher} reframes threat detection as a recommendation problem, training a GNN to predict whether information flows are benign or malicious.
More recent work has continued to diversify the detection landscape:
MirGuard~\cite{mirguard} combines noise injection with contrastive learning to improve robustness,
ProvADShield~\cite{provadshield} employs ensemble strategies across heterogeneous detector families,
MGDA~\cite{mgda} leverages graph reconstruction losses to surface anomalous substructures, and
Slot~\cite{slot} applies graph-level reinforcement learning for APT campaign detection.
The breadth of these approaches---spanning histograms, path embeddings, GNNs, and ensembles---motivates a detector-agnostic evasion strategy.
\textsc{SafeMimic} is designed and evaluated against both feature-based and GNN-based detectors, making no assumptions about the specific detection architecture.

\paragraph{Adversarial Attacks on Malware Detectors.}
A growing body of work studies adversarial manipulations that preserve malicious functionality while evading ML classifiers, though each is tailored to a specific artifact type.
HRAT~\cite{hrat} defines four deterministic atomic operations on function call graphs (FCGs) with predictable structural effects, enabling targeted graph perturbations.
BagAmmo~\cite{bagammo} injects dead code into Android FCGs via try-catch traps, ensuring that adversarial modifications never execute at runtime.
However, this paradigm does not transfer to provenance graphs, where every operation executes with real system effects and there is no notion of ``dead'' operations.
MalGuise~\cite{malguise} constructs a static safe transformation space for PE binaries using semantic NOPs and call-based redividing, demonstrating that safety-first design can achieve high evasion rates.
\textsc{SafeMimic} extends this safe-space-first paradigm to provenance graphs, but requires \emph{dynamic} safety reasoning---system state evolves with each injected operation, unlike MalGuise's static CPU semantics.
EvadeDroid~\cite{evadedroid} exploits component isolation in Android apps, where newly added components do not share state with existing ones, a property analogous to \textsc{SafeMimic}'s constraint $S_{\mathit{effect}} \cap S_{\mathit{dep}} = \emptyset$ but specific to the Android domain.
AdvDroidZero~\cite{advdroiedzero} uses UCB-based exploration to select perturbations without gradient access; such bandit strategies could complement \textsc{SafeMimic}'s greedy selection in future work.
ProvNinja~\cite{provninja} is the closest prior work, introducing regularity-based camouflage gadgets for provenance graph evasion.
However, ProvNinja lacks a safety mechanism: it selects operations solely by frequency in benign traces, which can create new frequency anomalies, corrupt attack dependencies, or trigger unintended system side effects---failure modes we formalize and empirically demonstrate in \S\ref{sec:evaluation}.

\paragraph{Adversarial Example Defenses.}
Defenses against adversarial examples provide additional context for \textsc{SafeMimic}'s design choices.
HagDe~\cite{hagde} detects adversarial examples by measuring decision boundary instability: inputs that lie close to the classifier's decision boundary exhibit inconsistent predictions under small perturbations.
This insight motivates \textsc{SafeMimic}'s strategy of gradual distribution shift rather than aggressive perturbation.
By matching the benign operation distribution via KL-divergence minimization, \textsc{SafeMimic} steers the provenance graph toward the interior of the benign region rather than toward boundary-adjacent areas that adversarial example detectors are designed to flag.
This distributional approach provides a degree of natural robustness against boundary-based defenses, though we discuss its limitations in \S\ref{sec:discussion}.
